\section{Qualitative affine-forest bounds for the sharp bootstrap}
\label{sec:path-qualitative-affine-bootstrap}

This section uses one qualitative marked-forest input.
It does not use a quantitative rate of convergence to one half.
All logarithms have base two.
A normalized $k$-pole is connected, has $k$ vertices of degree two, and has degree three at every other vertex.
Ranks are cycle-space dimensions, including for disconnected graphs.

\paragraph{Established forest input.}
Use Theorem~\ref{thm:sharp-qualitative-forest-input}, and increase its absolute rank threshold $R_*$ to at least $2^{100}$.
Its fixed-error conclusion is the following: if a connected graph has rank $\rho\ge R_*$ and a marked cycle $W$ with
\[
 |E(W)|\le\frac{\rho}{(\log\rho)^2},
\]
then
\begin{equation}\label{eq:path-qualitative-forest-input}
 \Pr(\xi\setminus E(W)\text{ is a forest})\le\frac{11}{20}.
\end{equation}
The input for simple graphs also applies to the localized multigraphs below.
Subdivide edges outside $W$ to remove loops and parallel edges; this preserves rank, the marked cycle, and the outside-forest probability.
Only this one-cycle input is used.

\begin{lemma}[Localization with a connected exterior]
\label{lem:path-qualitative-localization}
Let $G$ be connected, and let $R$ be a nonempty proper vertex set whose complement is connected.
For each edge of $\delta_G R$, add a distinct new vertex and an edge joining it to the original endpoint in $R$.
Join these new vertices in a cycle, padding to length three if necessary.
Call the new cycle $W_R$, its resulting graph $G_R$, and write $d=|\delta_G R|$.
Then $G_R$ is connected and
\[
 \beta(G_R)=\beta(G[R])+d-\operatorname{comp}(G[R])+1.
\]
Uniform even subgraphs have the same restriction law on $E(G[R])$ in $G$ and $G_R$.
If every vertex of $R$ is cubic in $G$, then
\[
 \beta(G_R)=\frac{|R|+d}{2}+1.
\]
\end{lemma}

\begin{proof}
Every component of $G[R]$ has a boundary edge, so the new cycle connects them all.
Before padding, the construction adds $d$ vertices and $2d$ edges.
Its rank is therefore $|E(G[R])|+d-|R|+1$, which is the first formula.
Padding adds equally many edges and vertices.
The cubic degree identity $3|R|=2|E(G[R])|+d$ gives the second formula.

Contract the connected complement of $R$ in $G$ and, separately, contract $W_R$ in $G_R$.
The resulting graphs agree on the original interior edges and the boundary edges.
Restriction under either contraction maps the original cycle space onto the quotient cycle space.
To see surjectivity, complete each prescribed even boundary by a join in a spanning tree of the contracted connected set.
A surjective linear map sends the uniform law to the uniform law on its image.
This proves the equality of restriction laws.
\end{proof}

\begin{lemma}[A fixed-error layer-width estimate]
\label{lem:path-qualitative-marker-width}
There is an absolute constant $B$ with the following property.
Let $G$ be connected and subcubic, and let $P$ be a nonempty connected marked set containing every vertex of degree below three.
Suppose that $P$ meets each distance layer from a fixed root in at most $k$ vertices, where $k\ge1$.
If
\[
 \Pr(\xi_G-P\text{ is a forest})\ge\frac35,
\]
then each full distance layer has at most $B(k+1)^2$ vertices.
\end{lemma}

\begin{proof}
Put $\Phi(m)=m/(\log m)^2$ and
\[
 \Theta=\max\{2R_*,\,2^{100}(k+1)^2\}.
\]
We first exclude every set $R\subseteq V(G)\setminus P$ with
\[
 m:=|R|\ge\Theta,
 \qquad |\delta_G R|\le\Phi(m)/4.
\]
Let $C_0$ be the component of $G-R$ containing $P$, and set
$R^+=V(G)\setminus V(C_0)$.
It contains $R$, avoids $P$, and has connected complement.
Distinct components of $G-R$ have no edge between them, so
$\delta_G R^+\subseteq\delta_G R$.
Every vertex of $R^+$ is cubic.

Localize at $R^+$ by Lemma~\ref{lem:path-qualitative-localization}.
Writing $d^+=|\delta_G R^+|$, the localized rank is
\[
 \rho=\frac{|R^+|+d^+}{2}+1\ge m/2\ge R_*.
\]
The function $\Phi$ is increasing on $[8,\infty)$.
If $d^+\ge3$, the marked cycle has order
\[
 d^+\le\Phi(m)/4\le\Phi(m/2)\le\Phi(\rho).
\]
If $d^+<3$, its padded order is three, which also satisfies this bound because $\rho\ge2^{100}$.
The event of a forest outside $P$ implies the forest event on $G[R^+]$.
In the localized graph, deleting the marked cycle edges leaves the new spokes incident with leaves.
Those spokes cannot create a cycle, so the outside-edge forest event is exactly the forest event on $G[R^+]$.
Localization and \eqref{eq:path-qualitative-forest-input} therefore bound its probability by $11/20$, contradicting the assumed lower bound $3/5$.
Thus every such $R$ is excluded.

Write $m_j$ for the number of unmarked vertices in distance layer $j$.
For the union $R(i,i')$ of unmarked layers $i$ through $i'$, its boundary satisfies
\[
 |\delta_G R(i,i')|
 \le3m_{i-1}+3m_{i'+1}+3k(i'-i+3).
\]
Indeed, every edge joins the same or consecutive layers; the last term counts incidences to marked vertices in layers $i-1$ through $i'+1$.
Use zero masses outside the layer range.

If some $m_j\ge\Theta$, begin with that layer.
Add a neighboring layer whenever its mass exceeds $m/(48(\log m)^2)$, where $m$ is the current mass.
The process stops in the finite graph.
Let $\mu$ be its final mass and $L=\log\mu\ge100$.
Each step multiplies the current mass by more than $1+1/(48L^2)$.
There are therefore at most
\[
 \frac{\ln\mu}{\ln(1+1/(48L^2))}
 \le96L^2\ln\mu<67L^3
\]
steps.
At termination the two neighboring masses are at most $\mu/(48L^2)$.
Hence
\[
 |\delta_G R(i,i')|
 \le\frac{\mu}{8L^2}+3k(67L^3+3)
 \le\frac{\mu}{4L^2}.
\]
For the last comparison, use
\[
 24kL^2(67L^3+3)\le2000kL^5\le2^L.
\]
The first inequality holds for $L\ge100$.
The threshold on $\mu$ gives $k\le2^{L/2-50}$, and
$2000L^5\le2^{L/2+50}$ for $L\ge100$.
The latter inequality holds at $100$, with an increasing right-to-left ratio thereafter.
The stopped window is thus an excluded set, a contradiction.

Every layer has fewer than $\Theta$ unmarked vertices and at most $k$ marked vertices.
For example,
\[
 B=2R_*+2^{100}+1
\]
bounds their total by $B(k+1)^2$.
\end{proof}

\begin{lemma}[Translation by a breadth-first tree marker]
\label{lem:path-qualitative-affine-width}
Let $N$ be a normalized $k$-pole with $k\ge2$.
If some even port subset $T$ has affine forest proportion at least $3/5$, then every distance layer from any chosen port has at most $B(k+1)^2$ vertices.
\end{lemma}

\begin{proof}
Choose a breadth-first tree rooted at that port and take the union of its root-to-port paths.
Its vertex set $P$ is connected, contains every degree-two vertex, and meets each distance layer in at most $k$ vertices.
The union of the tree edges contains a $T$-join $F_0$.
Translation by $F_0$ is a bijection from the even space to the affine $T$-join space.
The translated and original states agree outside $P$.
Thus an affine forest gives a state that is a forest outside $P$, and the preceding lemma applies.
\end{proof}

\begin{lemma}[Weighted Moore bound after suppressing the ports]
\label{lem:path-qualitative-weighted-moore}
Let $N$ be a normalized $k$-pole of rank $h\ge2$ and girth at least $D$.
With $t=\lceil\log(8h)\rceil$, one has
\[
 D\le4t\left(1+\frac{k}{3(h-1)}\right).
\]
\end{lemma}

\begin{proof}
Suppress all degree-two vertices and give each resulting edge the length of the path that it replaces.
The kernel is a connected cubic multigraph with
$n_0=2h-2$ vertices and $m_0=3h-3$ edges.
Its positive edge lengths sum to $m_0+k$.
Loops and parallel edges are permitted, and weighted cycles lift to cycles of the same length in $N$.

The kernel has $2m_0\,2^{t-1}$ directed non-backtracking walks of $t$ edges.
Under the uniform law on these walks, each directed edge is uniform at every step.
The expected total weight is $t(1+k/m_0)$.
At least half the walks have weight at most twice that expectation, and their number is at least
\[
 m_0\,2^{t-1}\ge12h(h-1)>n_0^2.
\]
If a short walk repeats a vertex, its support contains a cycle of no greater weight.
Otherwise all these short walks are simple paths, and two have the same ordered endpoints.
Their union contains a cycle of weight at most the sum of their weights.
In either case, a cycle has weight at most $4t(1+k/m_0)$.
This proves the bound.
\end{proof}

\begin{lemma}[The fixed-ratio affine input]
\label{lem:path-qualitative-fixed-ratio-core}
There is an absolute integer $D_*\ge3$ such that every normalized $k$-pole with
\[
 h\ge256,\qquad k\le4D,
 \qquad\operatorname{girth}(N)\ge D\ge D_*
\]
has affine forest proportion at most $3/5$ for every even port subset.
The assertion includes $k=0,1$.
\end{lemma}

\begin{proof}
With at most one port, the only even port subset is empty, and its forest proportion is $2^{-h}$.
Assume $k\ge2$ and put $t=\lceil\log(8h)\rceil$.
For $h\ge256$,
\[
 h-1>\frac{32}{3}\,t.
\]
Indeed, $t\le\log h+4$, the ratio $(h-1)/(\log h+4)$ increases for $h\ge2$, and at $h=256$ it exceeds $32/3$.
The weighted Moore bound therefore gives
\[
 D\le4t+\frac{16tD}{3(h-1)}<4t+\frac D2.
\]
Thus
\[
 h\ge2^{D/8-4},\qquad n=|V(N)|\ge h.
\]

Put $s=\lfloor D/16\rfloor$.
A graph of maximum degree three has at most $3\cdot2^s-2$ vertices in each radius-$s$ ball.
The union of these balls around the ports has at most
\[
 3k2^s\le12D2^{D/16}<2^{D/8-4}\le n
\]
vertices for sufficiently large $D$.
Choose a vertex outside that union.
Its radius-$s$ ball contains no port.
Since $2s<D$, its cubic breadth-first expansion has at least
$2^s\ge2^{D/16-1}$ vertices.
This ball meets at most $2s+1\le D$ distance layers from any fixed root.
Some layer therefore has at least $2^{D/16-1}/D$ vertices.

If any affine forest proportion were greater than $3/5$, Lemma~\ref{lem:path-qualitative-affine-width} would bound that layer by
$B(k+1)^2\le B(4D+1)^2$.
Choose $D_*$ large enough that the preceding ball comparisons hold and
\[
 \frac{2^{D/16-1}}D>B(4D+1)^2
 \qquad(D\ge D_*).
\]
The exponential term makes such a fixed threshold possible.
This contradiction proves the lemma.
\end{proof}

\begin{lemma}[An affine bound at the boundary scale]
\label{lem:path-clean-core-affine-forest}
There is an absolute constant $C_0$ such that every normalized $k$-pole with
\[
 h\ge L\ge4,
 \qquad k\le2L,
 \qquad\operatorname{girth}(N)\ge D\ge C_0\log(L+2)
\]
has affine forest proportion at most $3/5$ for every even port subset.
The assertion includes $k=0,1$.
\end{lemma}

\begin{proof}
The cases $k\le1$ again have forest proportion $2^{-h}$.
For $k\ge2$, one has
\[
 1+\frac{k}{3(h-1)}\le1+\frac{2h}{3(h-1)}<2.
\]
The weighted Moore bound gives
\[
 D<8\lceil\log(8h)\rceil\le8\log h+32,
 \qquad h\ge2^{D/8-4}.
\]
Choose $s=\lfloor D/16\rfloor$.
The radius-$s$ port balls contain at most $6L2^{D/16}$ vertices.
For a sufficiently large absolute $C_0$, this is less than
$n\ge h\ge2^{D/8-4}$ whenever $D\ge C_0\log(L+2)$.
The resulting port-free cubic ball gives a layer of size at least
$2^{D/16-1}/D$.

An affine forest proportion greater than $3/5$ would instead bound the layer by
$B(k+1)^2\le B(2L+1)^2$.
Increase $C_0$ so that the lower bound exceeds this quantity uniformly for $L\ge4$.
For explicit uniformity, the function $2^{D/16}/D$ increases for $D\ge32$.
At $D=C_0\log(L+2)$, its ratio to $(L+2)^2$ is
\[
 \frac{(L+2)^{C_0/16-2}}{C_0\log(L+2)}.
\]
This tends uniformly to infinity as $C_0\to\infty$, for $L\ge4$.
The same choice ensures the port-ball comparison.
The resulting contradiction proves the lemma.
\end{proof}

\begin{lemma}[Extracting a core from a large block]
\label{lem:path-clean-block-core}
Let $G$ be connected and cubic, and let $R$ be a nonempty proper vertex set with
\[
 |R|\ge8L^2,\qquad |\delta_G R|\le2L,\qquad L\ge4.
\]
Some connected component of $G[R]$ has a connected two-core of rank at least $L$, with at most $2L$ degree-two vertices and degree three elsewhere.
\end{lemma}

\begin{proof}
Write $n=|R|$, $d=|\delta_G R|$, and $c=\operatorname{comp}(G[R])$.
Each component has a boundary edge, so $1\le c\le d\le2L$.
The cubic degree identity gives
\[
 \beta(G[R])=\frac{n-d}{2}+c.
\]
Thus the average component rank is at least
\[
 \frac{n-d}{2d}+1
 \ge\frac{8L^2-2L}{4L}+1>L.
\]
Take a component of rank at least $L$ and prune its degree-one vertices.
The nonempty two-core remains connected and keeps the same rank.
At each step, the removed leaf has two incidences outside the current set and one inside it, so the boundary decreases by one.
The core boundary is therefore at most $2L$.
Its vertices have internal degree two or three, and each degree-two vertex has one boundary incidence.
This proves the port bound.
\end{proof}

\begin{proposition}[Two clean blocks]
\label{prop:path-two-clean-blocks}
Let $G$ be connected and cubic, and let $P$ be any marked set.
Suppose two vertex-disjoint sets $R_1,R_2\subseteq V(G)\setminus P$ satisfy
\[
 |R_i|\ge8L^2,\qquad |\delta_G R_i|\le2L,
 \qquad\operatorname{girth}(G[R_i])\ge D,
\]
where $L\ge4$ and $D\ge C_0\log(L+2)$.
Then
\[
 \Pr(\xi_G-P\text{ is a forest})\le\frac9{25}<\frac12.
\]
\end{proposition}

\begin{proof}
Extract cores $K_1,K_2$ by Lemma~\ref{lem:path-clean-block-core}.
Each remains disjoint from $P$, has rank at least $L$, at most $2L$ ports, and girth at least $D$.
Condition on every edge outside $E(K_1)\cup E(K_2)$.
For each feasible conditioning, the remaining states are independent and uniform in two affine parity classes.
The fiber is the product of their internal cycle spaces, translated by any compatible pair of states.
Only degree-two vertices can have odd prescribed parity, and each prescribed odd set has even size.

Lemma~\ref{lem:path-clean-core-affine-forest} bounds each conditional forest probability by $3/5$.
The forest event outside $P$ requires both core states to be forests.
Its conditional probability is therefore at most $9/25$.
Average over the conditioned edges.
Edges joining the two cores are already conditioned, so nonadjacency is not needed.
\end{proof}
